\documentclass[11pt]{article}
\usepackage[margin=1in]{geometry}
\usepackage{amsmath,amssymb,amsthm}
\usepackage{graphicx}
\usepackage{microtype}
\usepackage[hidelinks]{hyperref}

\newtheorem{theorem}{Theorem}
\newtheorem{corollary}[theorem]{Corollary}
\newcommand{\N}{\mathbb{N}}
\newcommand{\id}{\mathrm{id}}

\title{The Markov Group Conjecture for Triangular Countable-State Processes}
\author{Candidate partial result for arXiv:2004.05995}
\date{11 August 2026}

\begin{document}
\maketitle

\noindent\textbf{Status.} Candidate partial result, likely valid. This note
proves the triangular, or one-way acyclic, subcase and does not resolve the
full Markov group conjecture.

\section{Source conjecture}

Glueck restates the Markov group conjecture as Conjecture 1.1 of
\cite{glueck2020}, PDF page 1:

\begin{quote}
Let $T=(T_t)_{t\in[0,\infty)}$ be a Markovian $C_0$-semigroup on $\ell^1$
and assume that $T_1:\ell^1\to\ell^1$ is bijective (i.e., $T$ extends to a
$C_0$-group). Then $T$ has bounded generator.
\end{quote}

Here Markovian means that every $T_t$ is positive and norm-preserving on the
positive cone. Figure~\ref{fig:conjecture} is a crop of the source page.

\begin{figure}[ht]
  \centering
  \includegraphics[width=\textwidth]{figures/open_problem_crop.png}
  \caption{Conjecture 1.1 in \cite[p.~1]{glueck2020}.}
  \label{fig:conjecture}
\end{figure}

The source proves the conjecture for direct sums of finite-dimensional
processes, with a dimension-free bound, but explicitly leaves open the passage
from those decoupled systems to a general process.

\section{Triangular Markov semigroups}

Let $(e_j)_{j\in\N}$ be the canonical basis of $\ell^1(\N)$. A Markov
operator $T$ has a nonnegative column-stochastic matrix
\[
  p_{ij}=\langle Te_j,e_i^*\rangle,
  \qquad \sum_i p_{ij}=1.
\]
We call a Markov semigroup $(T_t)$ \emph{triangular} if, for every $t\ge0$,
\begin{equation}\label{eq:triangular}
  p_{ij}(t)=0 \qquad\text{whenever }i<j.
\end{equation}
Thus mass can move from state $j$ only to states $i\ge j$. A simultaneous
permutation of the states merely replaces the displayed order by another
fixed enumeration.

This class includes pure-birth and more general one-way processes with jumps
of arbitrary length. It need not decompose into finite invariant blocks.

\section{Idea of the proof}

For a triangular process, a path that starts at state $j$ and is at $j$ at a
later time could not have visited a larger state and returned. Algebraically,
the semigroup product has only one contribution to its $(j,j)$ entry. Hence
the diagonal probabilities are exponentials $e^{-q_jt}$.

The inverse bound at time one controls these exponentials. Indeed, the first
$j$ coordinate functionals form a finite-dimensional invariant subspace of
$T_1^*$, and its triangular matrix has $p_{jj}(1)$ as an eigenvalue. Thus
$p_{jj}(1)\ge\|T_1^{-1}\|^{-1}$ uniformly in $j$. This bounds all $q_j$.
Finally, the distance between a probability column and its starting basis
vector is exactly twice the mass that left the diagonal. Taking the supremum
over columns gives norm continuity at zero.

\section{Partial theorem and proof}

\begin{theorem}\label{thm:main}
Let $(T_t)_{t\ge0}$ be a triangular Markovian $C_0$-semigroup on
$\ell^1(\N)$. Suppose that $T_1$ is bijective, and put
\[
  M=\|T_1^{-1}\|.
\]
Then, for every $t\ge0$,
\begin{equation}\label{eq:normbound}
  \|T_t-\id\|\le 2(1-M^{-t}).
\end{equation}
In particular the generator $A$ of $(T_t)$ is bounded and
\begin{equation}\label{eq:generatorbound}
  \|A\|\le2\log M.
\end{equation}
\end{theorem}

\begin{proof}
Write $p_{ij}(t)=\langle T_te_j,e_i^*\rangle$. Positivity and preservation of
the $\ell^1$ norm on the positive cone give
\[
  p_{ij}(t)\ge0,
  \qquad \sum_i p_{ij}(t)=1.
\]
Fix $j$. By the semigroup law,
\[
  p_{jj}(t+s)=\sum_{k\in\N}p_{jk}(t)p_{kj}(s).
\]
Triangularity says that a nonzero factor $p_{kj}(s)$ requires $k\ge j$,
whereas a nonzero factor $p_{jk}(t)$ requires $j\ge k$. Thus only $k=j$
contributes, and
\begin{equation}\label{eq:multiplicative}
  p_{jj}(t+s)=p_{jj}(t)p_{jj}(s).
\end{equation}
Strong continuity gives $p_{jj}(0)=1$ and continuity in $t$. Consequently
there is a number $q_j\ge0$ such that
\begin{equation}\label{eq:diagonalexp}
  p_{jj}(t)=e^{-q_jt}\qquad(t\ge0).
\end{equation}

We next use bijectivity at time one. For $j\in\N$, let
\[
  F_j=\operatorname{span}\{e_1^*,\ldots,e_j^*\}\subset\ell^\infty.
\]
The subspace $F_j$ is invariant under $T_1^*$. Indeed, if $f\in F_j$ and
$k>j$, then
\[
  (T_1^*f)_k=\sum_i p_{ik}(1)f_i=0,
\]
because triangularity forces every potentially nonzero term to have
$i\ge k>j$, while $f_i=0$ there. In the displayed ordered basis, the
finite-dimensional restriction $T_1^*|_{F_j}$ is triangular and has diagonal
entries $p_{11}(1),\ldots,p_{jj}(1)$. Therefore $p_{jj}(1)$ is an eigenvalue
of $T_1^*$: choose a corresponding nonzero eigenvector $v\in F_j$.

Since $T_1^*$ is invertible and
$\|(T_1^*)^{-1}\|=\|T_1^{-1}\|=M$, the eigenvalue equation gives
\[
  \|v\|=\|(T_1^*)^{-1}p_{jj}(1)v\|
  \le M p_{jj}(1)\|v\|.
\]
Hence $p_{jj}(1)\ge M^{-1}$. By \eqref{eq:diagonalexp},
\begin{equation}\label{eq:qbound}
  q_j\le\log M\qquad(j\in\N).
\end{equation}

For each stochastic column, positivity and the column-sum identity yield
\begin{align*}
  \|(T_t-\id)e_j\|_1
  &=|p_{jj}(t)-1|+\sum_{i\ne j}p_{ij}(t)\\
  &=2(1-p_{jj}(t)).
\end{align*}
The norm of a matrix operator on $\ell^1$ is the supremum of the $\ell^1$
norms of its columns. Combining this fact with \eqref{eq:diagonalexp} and
\eqref{eq:qbound} gives
\[
  \|T_t-\id\|
  =\sup_j2(1-e^{-q_jt})
  \le2(1-e^{-t\log M})
  =2(1-M^{-t}),
\]
which is \eqref{eq:normbound}.

In particular $\|T_t-\id\|\to0$ as $t\downarrow0$. A $C_0$-semigroup that
is norm-continuous at zero has a bounded generator, and
\[
  A=\lim_{t\downarrow0}\frac{T_t-\id}{t}
\]
in operator norm \cite[Chapter~I]{engelnagel}. Dividing
\eqref{eq:normbound} by $t$ and passing to the limit proves
\eqref{eq:generatorbound}.
\end{proof}

\begin{corollary}
The Markov group conjecture is true for every countable-state process whose
transition semigroup becomes triangular after one fixed relabeling of the
states. If $\|T_1^{-1}\|=1$, then $T_t=\id$ for every $t\ge0$.
\end{corollary}

\begin{proof}
The first assertion is Theorem~\ref{thm:main}. For the second,
\eqref{eq:normbound} has right-hand side zero.
\end{proof}

\section{Verification notes}

\begin{itemize}
  \item No domain assumption on the basis vectors is made. The proof first
  obtains an operator-norm estimate and only then invokes the bounded-generator
  criterion.
  \item The eigenvalue argument occurs inside the genuinely invariant,
  finite-dimensional subspace $F_j$. It does not rely on a general claim that
  every diagonal entry of an infinite triangular operator lies in its
  spectrum.
  \item The triangular orientation is consistent throughout:
  $p_{ij}=0$ for $i<j$ means transitions $j\to i$ only with $i\ge j$.
  \item The exact column identity uses both positivity and conservation of
  mass. It would fail for merely substochastic semigroups.
  \item No numerical or symbolic computation is used.
\end{itemize}

\noindent\textbf{Verifier verdict:} likely valid. The most useful human
checks are the invariance of $F_j$ under $T_1^*$ and the equality
$\|T_t-\id\|=\sup_j2(1-p_{jj}(t))$.

\section{Scope and novelty status}

Cycles are the precise limitation: without triangularity, the $(j,j)$ entry
of $T_{t+s}$ includes paths that leave $j$ and later return, so
\eqref{eq:multiplicative} fails. The proof gives no control in that setting,
and the full Markov group conjecture remains open here.

On 11 August 2026, the local lightweight indexes and parsed arXiv corpus were
searched for the exact id, conjecture, and triangular, acyclic, and pure-birth
variants. Bounded official-arXiv searches used the same terms and found no
exact statement of this quantitative subcase. Novelty confidence is modest:
the argument is elementary, and older Markov-chain literature outside arXiv
may contain a comparable observation.

\section{Human review recommendation}

Verify the two points in the verdict. If accepted, retain this as a
quantitative partial result for triangular countable-state processes. It must
not be described as settling the full Markov group conjecture.

\begin{thebibliography}{9}
\bibitem{glueck2020}
Jochen Glueck,
\emph{On the decoupled Markov group conjecture},
arXiv:2004.05995, 2020.
\url{https://arxiv.org/abs/2004.05995}

\bibitem{engelnagel}
Klaus-Jochen Engel and Rainer Nagel,
\emph{One-Parameter Semigroups for Linear Evolution Equations},
Springer, 2000.
\end{thebibliography}

\end{document}

